%=============================================================================!
% 3.E  EXERCISES                                                              !
%=============================================================================!
\unnumbered{Exercises}
\label{sec:en-exercises}

The exercises run from questions of understanding, through calculations,
to short derivations marked `show that'. Full solutions are in
\hyperref[sec:en-solutions]{`Solutions to the exercises'} at the end of
the book, and Notebook~03 checks every numerical answer.

\begin{exercise}[a puck on a string]
\label{ex:en-circle}
A puck on a smooth, level table is tied by a string to a peg and slides
round the peg in a circle at a constant speed. Does the pull of the string
do work on the puck? What does your answer say about the puck's kinetic
energy?
\end{exercise}

\begin{exercise}[a suitcase]
\label{ex:en-suitcase}
A suitcase is pulled \SI{25}{\metre} along a level floor by its handle,
with a force of \SI{60}{\newton} directed \ang{50} above the horizontal.
How much work does the pull do?
\end{exercise}

\begin{exercise}[stairs]
\label{ex:en-stairs}
A person of \SI{70}{\kilogram} climbs a staircase that rises
\SI{3}{\metre}, at a steady speed, in \SI{5}{\second}. How much work do
they do against their weight, and at what average power?
\end{exercise}

\begin{exercise}[braking]
\label{ex:en-braking}
A car of \SI{1200}{\kilogram} travels at \SI{50}{\kilo\metre\per\hour}.
Find its kinetic energy, and the distance in which a braking force of
\SI{6000}{\newton} stops it. How far does the same force take to stop it
from \SI{100}{\kilo\metre\per\hour}?
\end{exercise}

\begin{exercise}[the second stretch]
\label{ex:en-stretch}
How much work does it take to stretch a spring of
\SI{200}{\newton\per\metre} from its natural length to \SI{10}{\centi
\metre}, and how much more to stretch it from \SI{10}{\centi\metre} to
\SI{20}{\centi\metre}?
\end{exercise}

\begin{exercise}[a launcher]
\label{ex:en-launcher}
A spring of \SI{500}{\newton\per\metre}, compressed by
\SI{8}{\centi\metre}, launches a puck of \SI{0.1}{\kilogram} along a
smooth level table. Find the puck's speed once it leaves the spring, and
the height it reaches if it then slides up a smooth ramp.
\end{exercise}

\begin{exercise}[a pendulum]
\label{ex:en-pendulum}
A small bob hangs from a string \SI{1.2}{\metre} long and is released
from rest with the string at \ang{40} to the vertical. Explain why the
string does no work on the bob, and find the bob's speed at the bottom of
its swing.
\end{exercise}

\begin{exercise}[back on the track]
\label{ex:en-track}
The bead on the track of \cref{eq:en-track} passes the bottom of the deep
valley at \SI{3}{\metre\per\second}. What is the greatest height it
reaches? Can it cross the hill? Use \cref{fig:en-track} to estimate its
turning points, and Notebook~03 to find them.
\end{exercise}

\begin{exercise}[a double well]
\label{ex:en-double}
A body moves along a line with potential energy $U(x) = U_0\,[(x/a)^2 -
1]^2$, where $U_0$ and $a$ are positive constants. Show that it has
equilibria at $x = 0$ and $x = \pm a$, and decide from $U''$ which are
stable. For $E = U_0/2$, show that the turning points are at $x = \pm
a\sqrt{1 \pm 1/\sqrt2}$, and describe the allowed regions.
\end{exercise}

\begin{exercise}[on ice]
\label{ex:en-ice}
A puck slides onto rough ice at \SI{3}{\metre\per\second}, with
$\mu_{\mathrm{k}} = 0.05$. How far does it slide before it stops?
\end{exercise}

\begin{exercise}[a rough ramp]
\label{ex:en-ramp}
A block slides from rest down a straight ramp inclined at \ang{25},
starting \SI{2}{\metre} above the bottom, with $\mu_{\mathrm{k}} = 0.15$.
Find its speed at the bottom.
\end{exercise}

\begin{exercise}[other carts]
\label{ex:en-carts}
Carts of \SI{1}{\kilogram} and \SI{2}{\kilogram}, at rest on a smooth
track with the lighter on the left, are pushed apart by a spring of
\SI{200}{\newton\per\metre} compressed by \SI{5}{\centi\metre}. Find
their velocities, with $x$ to the right, once the spring has fallen away.
\end{exercise}

\begin{exercise}[coupling]
\label{ex:en-coupling}
In \cref{ex:ne-coupling} a cart of \SI{2}{\kilogram} moving at
\SI{3}{\metre\per\second} couples to a cart of \SI{1}{\kilogram} at rest,
and the pair moves on at \SI{2}{\metre\per\second}. Compare the kinetic
energies before and after. Where has the difference gone, and why does
momentum not show the same loss?
\end{exercise}

\begin{exercise}[two routes]
\label{ex:en-routes}
A force in the plane is $\vb{F} = a\,(y, x)$, with $a =
\SI{0.5}{\electronvolt\per\angstrom\squared}$. Compute its work from the
origin to $(2, 1)$\,\si{\angstrom} along the straight line, and along
the route that goes first along the $x$ axis to $(2, 0)$ and then
straight up. Find a potential energy $U$ with $\vb{F} = -\grad U$.
\end{exercise}

\begin{exercise}[an elliptical well]
\label{ex:en-ellipse}
An atom in the plane has potential energy $U = \frac12 k\,(x^2 + 4y^2)$,
with $k = \SI{2}{\electronvolt\per\angstrom\squared}$. Find the force on
it at $(1, 1)$\,\si{\angstrom}. Show that the force is perpendicular to
the contour of $U$ through that point, and that it does not point at the
origin, where $U$ is least.
\end{exercise}

\begin{exercise}[the curl test]
\label{ex:en-curl}
Show that $\vb{F} = a\,(2xy,\; x^2 + z^2,\; 2yz)$, with $a$ in
\si{\electronvolt\per\angstrom\cubed}, has zero curl, and find the
potential energy $U$ with $U(\vb{0}) = 0$. Evaluate $U$ at $(1, 2,
3)$\,\si{\angstrom} for $a = \SI{1}{\electronvolt\per\angstrom\cubed}$.
\end{exercise}

\begin{exercise}[round a square]
\label{ex:en-square}
Compute the work of the swirl $\vb{F} = c\,(-y, x)$ once anticlockwise
around the square with corners $(0, 0)$, $(L, 0)$, $(L, L)$ and $(0, L)$,
side by side, and show that it is $2c$ times the area of the square.
\end{exercise}

\begin{exercise}[a field with a hole]
\label{ex:en-vortex}
In the plane with the origin removed, let $\vb{F} = b\,(-y, x)/(x^2 +
y^2)$, with $b$ in \si{\electronvolt}. Show that $\partial F_y/\partial x
= \partial F_x/\partial y$ wherever $\vb{F}$ is defined (differentiate
$(x^2 + y^2)\,q = 1$ by the product rule to find the partial derivatives
of $q = 1/(x^2 + y^2)$), and compute the
work once anticlockwise around the circle of radius $R$ about the origin.
Explain why a non-zero answer does not contradict the curl test, and find
the work around a circle that does not enclose the origin.
\end{exercise}

\begin{exercise}[the top site]
\label{ex:en-top}
For the model surface of \cref{eq:en-surface}, write $\vb{g}_1$,
$\vb{g}_2$ and $\vb{g}_3$ in components, and show by the scalar product in
components that at the top $\vb{r} = (a/2, a/(2\sqrt3))$ the three
products $\vb{g}_k\cdot\vb{r}$ are $2\pi/3$, $2\pi/3$ and $-4\pi/3$, and
hence that $U = 9U_{\mathrm b}/8$.
\end{exercise}

\begin{exercise}[pair forces]
\label{ex:en-pair}
Two atoms have potential energy $U = \varphi(r_{12})$, a function of their
distance $r_{12} = \lvert\vb{r}_2 - \vb{r}_1\rvert$ alone. Show that
$\partial r_{12}/\partial x_2 = (x_2 - x_1)/r_{12}$ (differentiate
$r_{12}^2$, a sum of three squares, with respect to $x_2$), and hence
that
$\vb{F}_2 = -\varphi'(r_{12})\,(\vb{r}_2 - \vb{r}_1)/r_{12}$ and
$\vb{F}_1 = -\vb{F}_2$: the forces are equal and opposite and act along
the line joining the atoms.
\end{exercise}

\begin{exercise}[stopping an oxygen atom]
\label{ex:en-oxygen}
An oxygen atom (\SI{15.999}{\amu}) moving at
\SI{0.01}{\angstrom\per\femto\second} enters a region where a constant
force of \SI{0.5}{\electronvolt\per\angstrom} opposes its motion. Use the
work-energy theorem to find how far it travels before it stops.
\end{exercise}

\begin{exercise}[a drag on an atom]
\label{ex:en-drag}
An atom moves with no potential energy under the drag $-\gamma m\vb{v}$.
Use \cref{eq:en-dEdt} to show that $K(t) = K(0)\,\mathrm{e}^{-2\gamma t}$.
For a lithium atom (\SI{6.94}{\amu}) starting at
\SI{0.02}{\angstrom\per\femto\second} with $\gamma =
\SI{0.005}{\per\femto\second}$, find $K$ after \SI{100}{\femto\second},
in electronvolts.
\end{exercise}
